\corrigesChapitre{7}{La matière dans l’Univers}
\par\noindent\begin{minipage}{\linewidth}
\begin{corrige}{1}{Classer des objets}
\begin{questions}
\item Soleil : étoile ; Terre : planète ; Lune : satellite naturel.
\item Oui, comme l’ensemble du Système solaire.
\end{questions}
\end{corrige}
\end{minipage}
\par\noindent\begin{minipage}{\linewidth}
\begin{corrige}{2}{Le bon emboîtement}
\begin{questions}
\item Terre, Système solaire, Voie lactée, Univers observable, en allant vers l’ensemble le plus vaste.
\item Non : elle rassemble notamment de très nombreuses étoiles, du gaz et des poussières.
\end{questions}
\end{corrige}
\end{minipage}
\par\noindent\begin{minipage}{\linewidth}
\begin{corrige}{3}{Une unité trompeuse}
\begin{questions}
\item Une distance.
\item Elle correspond au trajet de la lumière pendant une durée d’un an.
\end{questions}
\end{corrige}
\end{minipage}
\par\noindent\begin{minipage}{\linewidth}
\begin{corrige}{4}{L’ordre des planètes}
\begin{questions}
\item Vénus avant et Mars après.
\item Neptune.
\end{questions}
\end{corrige}
\end{minipage}
\par\noindent\begin{minipage}{\linewidth}
\begin{corrige}{5}{Lire une échelle}
\begin{questions}
\item $50/10=5\unite{cm}$.
\item $150/10=15\unite{cm}$.
\end{questions}
\pcFigureEx{figures/cor-echelle-astronomie.pdf}{Schéma réduit : les longueurs indiquées sont à reporter sur la frise. L’échelle relie chaque centimètre à dix millions de kilomètres ; les diamètres ne sont pas à l’échelle.}
\end{corrige}
\end{minipage}
\par\noindent\begin{minipage}{\linewidth}
\begin{corrige}{6}{Comparer des distances}
\begin{questions}
\item $600/150=4$.
\item $600-150=450$ millions de kilomètres.
\end{questions}
\end{corrige}
\end{minipage}
\par\noindent\begin{minipage}{\linewidth}
\begin{corrige}{7}{Des milliards d’années}
\begin{questions}
\item $4\,600$ millions d’années.
\item C’est une échelle astronomique et géologique, très supérieure à une durée humaine.
\end{questions}
\end{corrige}
\end{minipage}
\par\noindent\begin{minipage}{\linewidth}
\begin{corrige}{8}{Les mêmes éléments}
\begin{questions}
\item Non : le même symbole H désigne le même élément, soumis aux mêmes lois.
\item Les milieux ont des histoires et compositions différentes ; mêmes éléments possibles ne signifie pas mêmes proportions.
\end{questions}
\end{corrige}
\end{minipage}
\par\noindent\begin{minipage}{\linewidth}
\begin{corrige}{9}{Les limites d’une affiche}
\begin{questions}
\item Non : elle peut montrer l’ordre ou les noms.
\item Préciser que les diamètres ne sont pas à l’échelle ; indiquer séparément le choix éventuel d’une échelle de distances.
\end{questions}
\end{corrige}
\end{minipage}
\par\noindent\begin{minipage}{\linewidth}
\begin{corrige}{10}{Passer à une autre échelle}
\begin{questions}
\item $750/25=30\unite{cm}$ du Soleil.
\item $8\times25=200$ millions de kilomètres.
\end{questions}
\pcFigureEx{figures/cor-frise-750.pdf}{On divise la distance réelle par la distance représentée par un centimètre.}
\end{corrige}
\end{minipage}
\par\noindent\begin{minipage}{\linewidth}
\begin{corrige}{11}{Une preuve d’autorité ?}
\begin{questions}
\item Arrondis différents ou distance instantanée comparée à une distance moyenne ; les orbites ne gardent pas toujours la même distance.
\item La définition de la grandeur, la date, l’unité et l’auteur de la source.
\end{questions}
\end{corrige}
\end{minipage}
\par\noindent\begin{minipage}{\linewidth}
\begin{corrige}{12}{Un métal « rare »}
\begin{questions}
\item Non : il peut être très dispersé et difficile à concentrer.
\item La teneur des gisements, l’accessibilité, l’énergie nécessaire, les déchets et les possibilités de recyclage.
\end{questions}
\end{corrige}
\end{minipage}
\par\noindent\begin{minipage}{\linewidth}
\begin{corrige}{13}{Une autre unité}
\begin{questions}
\item $450/150=3\unite{ua}$.
\item $5\times150=750$ millions de kilomètres.
\end{questions}
\end{corrige}
\end{minipage}
\par\noindent\begin{minipage}{\linewidth}
\begin{corrige}{14}{Un alignement dessiné}
\begin{questions}
\item Non : il ne donne pas leurs positions orbitales à une date précise.
\item Il peut représenter l’ordre des distances moyennes et les noms, selon sa légende.
\end{questions}
\end{corrige}
\end{minipage}
\par\noindent\begin{minipage}{\linewidth}
\begin{corrige}{15}{Rédiger une histoire}
\begin{questions}
\item La matière terrestre provient de matière cosmique enrichie au cours de l’évolution stellaire avant la formation du Système solaire.
\item La Terre a hérité de nombreux éléments ; elle n’est pas l’origine unique de ces éléments.
\end{questions}
\end{corrige}
\end{minipage}
\par\noindent\begin{minipage}{\linewidth}
\phantomsection\label{cor-defi-7}\pcSolutionDefi{7}{La frise impossible ?}\begin{questions}
\item $9\unite m=900\unite{cm}$ ; on peut choisir $1\unite{cm}$ pour $5$ millions de kilomètres.
\item $150/5=30\unite{cm}$ du Soleil.
\item Les diamètres sont minuscules devant les distances orbitales. Il faut employer des symboles et annoncer explicitement la différence d’échelle.
\end{questions}
\pcFigureEx{figures/cor-frise-salle.pdf}{Une longueur totale de 900 cm représente 4 500 millions de kilomètres ; la Terre est à 30 cm du Soleil.}
\end{minipage}
\par\noindent\begin{minipage}{\linewidth}
\phantomsection\label{cor-activite-7}\pcSousTitre{Pistes pour l’investigation}
Calculer chaque position en divisant la distance réelle par la valeur représentée par une unité de longueur. Choisir la plus grande distance pour vérifier que toute la frise tient dans la salle. Mentionner séparément les diamètres hors échelle. Une agence spatiale ou un organisme scientifique identifié fournit une source de référence ; les valeurs arrondies ne sont pas des mesures instantanées.\par\end{minipage}
